% Generated draft: requires mathematical and source review.
\documentclass[11pt]{article}
\usepackage[margin=25mm]{geometry}
\usepackage{fontspec}
\setmainfont{Latin Modern Roman}
\usepackage{amsmath,amssymb,mathtools,unicode-math}
\setmathfont{Latin Modern Math}
\usepackage{enumitem,needspace,xurl,hyperref}
\setlength{\parindent}{0pt}
\setlength{\parskip}{5pt}
\setlength{\emergencystretch}{3em}
\newcommand{\N}{\mathbb{N}}
\newcommand{\Z}{\mathbb{Z}}
\newcommand{\Q}{\mathbb{Q}}
\newcommand{\R}{\mathbb{R}}
\newcommand{\C}{\mathbb{C}}
\newcommand{\F}{\mathbb{F}}
\newcommand{\A}{\mathbb{A}}
\newcommand{\PP}{\mathbb{P}}
\newcommand{\E}{\mathbb{E}}
\newcommand{\eps}{\varepsilon}
\newcommand{\dd}{\,\mathrm{d}}
\begin{document}
\section*{UnsolvedMath ID: 800098}
\textbf{Source problem code:} catalogue entry
\subsection{Definitions and mathematical statement}
The supplied record leaves the topology on the representation space and the definition of the Euler number implicit. The following is a conventional interpretation; these conventions require confirmation against the cited source.

Let \(G=\mathrm{Homeo}_0(\R^2)\) denote the identity component of the homeomorphism group of the plane, equipped with the compact-open topology. Give
\[
\mathrm{Hom}(\Z^2,G)
\cong \{(f,g)\in G\times G: f\circ g=g\circ f\}
\]
the subspace topology inherited from \(G\times G\), where the identification uses the ordered standard generators of \(\Z^2\).

To specify an Euler-number convention, let \(e\in H^2(BG;\Z)\) be the universal Euler class, normalized so that its restriction along the standard rotation subgroup \(\mathrm{SO}(2)\hookrightarrow G\) is the usual Euler class of the universal oriented real two-plane bundle. Here \(BG\) denotes the classifying space of the topological group \(G\). For a representation \(\rho:\Z^2\to G\), form the associated flat plane bundle
\[
E_\rho=(\R^2\times\R^2)/\Z^2\longrightarrow T^2=\R^2/\Z^2,
\qquad
m\cdot(x,y)=(x+m,\rho(m)(y)).
\]
Let \(c_\rho:T^2\to BG\) be its classifying map. Orient \(T^2\) using the ordered standard basis of \(\Z^2\), and define
\[
\operatorname{eu}(\rho)
=\bigl\langle c_\rho^*e,[T^2]\bigr\rangle\in\Z,
\qquad
\mathrm{Hom}_k(\Z^2,G)
=\{\rho\in\mathrm{Hom}(\Z^2,G):\operatorname{eu}(\rho)=k\},
\]
with the induced subspace topology.

Under these conventions, the question supplied as Question 39 is:
\[
\text{For every }k\in\Z,\quad
\mathrm{Hom}_k\bigl(\Z^2,\mathrm{Homeo}_0(\R^2)\bigr)
\text{ is connected?}
\]
Here connected means that the space admits no decomposition into two disjoint nonempty relatively open subsets. The question does not assert that every Euler-number level is nonempty. The supplied record labels the problem open; that label is not an independently verified mathematical assessment.
\subsection{Short English statement}
For each integer \(k\), consider all actions of \(\Z^2\) on the plane by homeomorphisms in the identity component whose Euler number is \(k\). Is this space connected? The draft uses the compact-open topology and an explicitly normalized Euler class; the supplied record does not specify these conventions.
\subsection{Sources}
\begin{itemize}
\item Local UnsolvedMath registry, permanent ID \(800098\), exact source problem code \texttt{catalogue entry}, title ``Foliations and Diffeomorphism Groups --- Question 39.'' The supplied record attributes the problem to Nariman and Yazdi, gives proposed year 2025, and labels its status \texttt{open}. Upstream revision: \texttt{d371ee917ac9647fbb06cd6fc529deec6f3dda82}.
\item Nariman and Yazdi, \emph{Problem List on Foliations and Diffeomorphism Groups} (2025), Question 39, \url{https://arxiv.org/abs/2503.18274}. The supplied provenance records access on 2026-10-16 and extraction by \texttt{arxiv\_latex}. This access information is supplied provenance, not a claim of independent source consultation or verification.
\end{itemize}
\end{document}
